% Gemeinsame Praeambel fuer alle drei Zusammenfassungen.
% Wird von main_performance.tex, main_docker.tex und main_power.tex
% eingebunden. Hier aendern wirkt auf alle drei Dokumente.
% Jede main-Datei definiert vorher \doctitle und \docsubject.

% ---------------------------------------------------------------------
% Pakete
% ---------------------------------------------------------------------
\usepackage[utf8]{inputenc}
\usepackage[T1]{fontenc}
\IfFileExists{ngerman.ldf}{\usepackage[ngerman]{babel}}{%
  \renewcommand{\contentsname}{Inhaltsverzeichnis}}
\usepackage{helvet}
\renewcommand{\familydefault}{\sfdefault}

\usepackage[a4paper,left=1.0cm,right=1.0cm,top=1.0cm,bottom=1.1cm,
            footskip=0.5cm]{geometry}
\usepackage{amsmath,amssymb}
\usepackage{graphicx}
\usepackage{multicol}
\usepackage{booktabs}
\usepackage{array}
\usepackage{tabularx}
\usepackage{environ}        % damit tabularx in \ztab funktioniert
\usepackage{enumitem}
\usepackage[dvipsnames]{xcolor}
\usepackage{titlesec}
\usepackage[most]{tcolorbox}
\usepackage{listings}
\usepackage{fancyhdr}
\usepackage{lastpage}
\usepackage[colorlinks=true,linkcolor=blue!60,urlcolor=blue!60]{hyperref}

\graphicspath{{figures/}}

% ---------------------------------------------------------------------
% Metadaten (hier anpassen)
% ---------------------------------------------------------------------
\newcommand{\docauthors}{Jan Vogel}
\newcommand{\docdate}{\today}
\newcommand{\docversion}{v3.0}

% ---------------------------------------------------------------------
% Layout: 2 Spalten, Abstände
% Platz knapp? Zuerst hier an den Rändern drehen, dann an der
% Schriftgrösse in \documentclass ganz oben.
% ---------------------------------------------------------------------
\setlength{\columnsep}{10pt}
\setlength{\columnseprule}{0.4pt}
\renewcommand{\columnseprulecolor}{\color{black!30}}
\setlength{\parindent}{0pt}
\setlength{\parskip}{0.5pt}
\setlength{\tabcolsep}{3pt}
\setcounter{secnumdepth}{2}
\setcounter{tocdepth}{1}

\setlist[itemize]{nosep,leftmargin=1.1em,topsep=1pt,label=\textcolor{accent}{\textbullet}}
\setlist[enumerate]{nosep,leftmargin=1.5em,topsep=1pt}
\setlist[description]{nosep,topsep=1pt,leftmargin=0pt,font=\normalfont\bfseries}

\setlength{\jot}{5pt}
\allowdisplaybreaks
\AtBeginDocument{%
  \setlength{\abovedisplayskip}{3pt}\setlength{\belowdisplayskip}{3pt}%
  \setlength{\abovedisplayshortskip}{3pt}\setlength{\belowdisplayshortskip}{3pt}}

% ---------------------------------------------------------------------
% Farbe: eine durchgehende Akzentfarbe fuer das ganze Dokument.
% Die Boxenfarben stehen fuer die ART des Inhalts (Falle, Rezept,
% Beispiel) -- nicht fuer eine Vorlesungswoche.
% ---------------------------------------------------------------------
\definecolor{accent}{HTML}{1F4E79}
\colorlet{accentbar}{accent!14}

% ---------------------------------------------------------------------
% Überschriften
% ---------------------------------------------------------------------
\titleformat{\section}[block]{\normalfont}{}{0pt}{\sectionbar}
\titleformat{name=\section,numberless}[block]{\normalfont}{}{0pt}{\sectionbarnonum}
\titleformat{\subsection}[block]{\normalfont}{}{0pt}{\subsectionbar}
\newcommand{\sectionbar}[1]{\colorbox{accentbar}{\parbox{\dimexpr\linewidth-2\fboxsep\relax}{\large\strut\textcolor{accent}{\textbf{\thesection\ #1}}}}}
\newcommand{\sectionbarnonum}[1]{\colorbox{accentbar}{\parbox{\dimexpr\linewidth-2\fboxsep\relax}{\large\strut\textcolor{accent}{\textbf{#1}}}}}
\newcommand{\subsectionbar}[1]{{\color{accent}\rule{2pt}{8pt}}\hskip3pt{\footnotesize\textbf{\thesubsection\ #1}}}

\titlespacing*{\section}{0pt}{3pt}{1pt}
\titlespacing*{\subsection}{0pt}{3pt}{0.5pt}

% ---------------------------------------------------------------------
% Kopf-/Fusszeile
% ---------------------------------------------------------------------
\pagestyle{fancy}
\fancyhf{}
\renewcommand{\headrulewidth}{0pt}
\renewcommand{\footrulewidth}{0.4pt}
\fancyfoot[L]{\tiny \docauthors\ --- EPP \docsubject\ --- \docversion}
\fancyfoot[R]{\tiny \thepage/ \pageref{LastPage}}

% ---------------------------------------------------------------------
% Eigene Umgebungen und Befehle
% ---------------------------------------------------------------------
\newcommand{\R}[1]{\texttt{#1}}          % Register / Code / Befehl
\newcommand{\T}[1]{\textbf{#1}}          % Schluesselbegriff
\newcommand{\ra}{\,$\rightarrow$\,}
\newcommand{\SW}[1]{\textcolor{accent}{\textbf{#1}}\hskip0.4em}  % Stichwort am Zeilenanfang
\newcommand{\merkzeile}{\par\textcolor{accent}{$\blacktriangleright$}\ } % Merksatz-Zeile

% Kompakte Tabelle ueber die Spaltenbreite. Erste Zeile = Kopfzeile.
%   \begin{ztab}{@{}Xr@{}} Kopf & Kopf \\ \midrule ... \end{ztab}
% (\NewEnviron statt \newenvironment, sonst findet tabularx sein Ende nicht)
\NewEnviron{ztab}[1]{%
  \par\footnotesize
  \begin{tabularx}{\linewidth}{#1}\toprule\BODY\bottomrule\end{tabularx}%
  \par\vspace{1pt}\normalsize}

% Boxen -- die Farbe sagt, WAS fuer ein Inhalt drinsteht
\tcbset{zfbox/.style={enhanced,arc=1.5mm,boxrule=0.6pt,
  left=3pt,right=3pt,top=2pt,bottom=2pt,
  fonttitle=\sffamily\bfseries\footnotesize,fontupper=\footnotesize,
  before skip=2pt,after skip=2pt}}

\newtcolorbox{formel}[1][]{zfbox,colback=accent!4,colframe=accent,
  coltitle=white,colbacktitle=accent,title={#1}}         % Formel/Kernaussage
\newtcolorbox{falle}[1][Falle]{zfbox,colback=red!4,colframe=red!65,
  coltitle=white,colbacktitle=red!65,title={#1}}         % typischer Fehler
\newtcolorbox{rezept}[1][Rezept]{zfbox,colback=green!4,colframe=green!40!black,
  coltitle=white,colbacktitle=green!40!black,title={#1}} % Schritt fuer Schritt
\newtcolorbox{bsp}[1][Beispiel]{zfbox,colback=yellow!8,colframe=black!55,
  coltitle=white,colbacktitle=black!55,title={#1}}       % durchgerechnet
\newtcolorbox{achtung}[1][Achtung]{zfbox,colback=red!4,colframe=red!65,
  coltitle=white,colbacktitle=red!65,title={#1}}

% Bild in Spaltenbreite (Anteil von \linewidth)
\newcommand{\fig}[2][0.9]{\par\vspace{2pt}{\centering\includegraphics[width=#1\linewidth]{#2}\par}\vspace{1pt}}

% Code
\lstset{
  basicstyle=\ttfamily\scriptsize,
  columns=fullflexible, keepspaces=true,
  aboveskip=2pt, belowskip=2pt, xleftmargin=3pt,
  commentstyle=\color{black!55}\itshape,
  morecomment=[l]{\#}, morecomment=[l]{;},
  literate={ä}{{\"a}}1 {ö}{{\"o}}1 {ü}{{\"u}}1 {ß}{{\ss}}1
}

% =====================================================================
